\section{Conclusion}
\label{sec:conclusion}
% ============================================================

We have presented a comprehensive yield system for cross-chain bridge tokens that eliminates the \$675M+ annual capital inefficiency of idle bridge TVL. The system's formal properties---share price monotonicity, automatic undercollateralization pause, tier-based diversification, and exit guarantees---provide the safety invariants required for production bridge operation. Production hardening includes \texttt{VIRTUAL\_SHARES} $= \texttt{VIRTUAL\_ASSETS} = 10^8$ for first-depositor inflation attack protection, per-operation \texttt{forceApprove} (no infinite approvals), and 54 security tests including fuzz (256 runs) and invariant (128K calls) suites.

The Bitcoin-native yield composition through Babylon, Lombard, SolvBTC, and CoreDAO represents a qualitative advance: native BTC earning yield without custodial wrapping, with the same FROST threshold key securing both bridge custody and staking positions.

The xLUX flywheel, where fees from bridge, DEX, lending, perpetuals, and NFT AMM compound into staking yield, creates a self-reinforcing economic equilibrium that scales with ecosystem adoption. As cross-chain volume grows, yield-bearing bridge tokens become strictly dominant over conventional wrapped tokens---offering higher capital efficiency with equivalent or superior security guarantees.

% ============================================================
% References
% ============================================================
\begin{thebibliography}{99}

\bibitem{defillama2025} DefiLlama, ``Bridge TVL Rankings,'' \url{https://defillama.com/bridges}, accessed April 2025.

\bibitem{luxteleport} Z.~Kelling, ``Lux Teleport: A Sovereign Omnichain ,'' Lux Industries, Inc., 2025.

\bibitem{frost} C.~Komlo and I.~Goldberg, ``FROST: Flexible Round-Optimized Schnorr Threshold Signatures,'' in \emph{Selected Areas in Cryptography (SAC)}, 2020.

\bibitem{cggmp21} R.~Canetti, R.~Gennaro, S.~Goldfeder, N.~Makriyannis, and U.~Peled, ``UC Non-Interactive, Proactive, Threshold ECDSA with Identifiable Aborts,'' in \emph{ACM CCS}, 2020.

\bibitem{erc4626} J.~Bebis \emph{et al.}, ``EIP-4626: Tokenized Vault Standard,'' Ethereum Improvement Proposal, 2022.

\bibitem{lidofinance} Lido DAO, ``Lido:  for Ethereum,'' \url{https://lido.fi}, 2020.

\bibitem{aavev3} Aave Companies, ``Aave V3 Technical Paper,'' 2022.

\bibitem{eigenlayer2023} S.~Kannan \emph{et al.}, ``EigenLayer: The Restaking Collective,'' Whitepaper, 2023.

\bibitem{babylon2024} Babylon Labs, ``Bitcoin Staking: Unlocking 21M Bitcoins to Secure the Proof-of-Stake Economy,'' Whitepaper, 2024.

\bibitem{lombard2024} Lombard Finance, ``LBTC: Liquid Bitcoin Staking,'' \url{https://lombard.finance}, 2024.

\bibitem{solvbtc} Solv Protocol, ``SolvBTC: Bitcoin Yield Aggregation,'' \url{https://solv.finance}, 2024.

\bibitem{coredao} Core Foundation, ``CoreDAO: Satoshi Plus Consensus,'' Whitepaper, 2023.

\bibitem{pendle2023} Pendle Finance, ``Pendle: Yield Tokenization Protocol,'' \url{https://pendle.finance}, 2023.

\bibitem{across2023} Across Protocol, ``Across: Optimistic Cross-Chain Bridge,'' Whitepaper, 2023.

\bibitem{wormhole} Wormhole Foundation, ``Wormhole: A Generic Message Passing Protocol,'' Whitepaper, 2022.

\bibitem{wbtc} BitGo, ``Wrapped Bitcoin (WBTC),'' \url{https://wbtc.network}, 2019.

\bibitem{rekt} Rekt News, ``Leaderboard,'' \url{https://rekt.news/leaderboard/}, 2024.

\bibitem{markowitz1952} H.~Markowitz, ``Portfolio Selection,'' \emph{The Journal of Finance}, vol.~7, no.~1, pp.~77--91, 1952.

\end{thebibliography}

